\documentclass[a4paper,11pt]{article}
% Paquetes necesarios
\usepackage[utf8]{inputenc}  % Para caracteres en español
\usepackage{amsmath, amssymb} % Para símbolos matemáticos
\usepackage{fancyhdr} % Para encabezado y pie de página
\usepackage{geometry} % Para márgenes
\usepackage{enumitem} % Para modificar listas
\usepackage{multicol} % Para múltiples columnas

% Configuración de márgenes
\geometry{left=2cm, right=2cm, top=2cm, bottom=2cm}
\setlength{\headheight}{14pt}

% Configuración del encabezado y pie de página
\pagestyle{fancy}
\lfoot{Ing. Luis A. Molina}  % Pie de página a la izquierda
\cfoot{MAT207}  % Centro vacío
\rfoot{\thepage} % Número de página a la derecha

% Inicio del documento
\begin{document}

% Título de la práctica
\begin{center}
    \Large\textbf{Práctica 07 - Ecuaciones Diferenciales Lineales de Orden Superior}\\[0.5cm]  
\end{center}
  \textbf{Nombre:} \rule{9cm}{0.4pt}  \textbf{Fecha:} \rule{4.5cm}{0.4pt}

% Introducción
\section*{Introducción}
En esta práctica, se resolverán ecuaciones diferenciales lineales homogéneas de orden $n$ con coeficientes constantes reales: el cálculo de la solución general $y_h(x)$ mediante la ecuación característica, la resolución de problemas de valor inicial (PVI), la reconstrucción de operadores diferenciales a partir de conjuntos fundamentales de soluciones, y problemas de deducción conceptual sobre raíces complejas y orden mínimo.

% Ejercicios
\section*{Ejercicios}

\subsection*{A) Solución General de Ecuaciones Homogéneas de Orden $n$}
Halla la solución general $y_h(x)$ para cada una de las siguientes ecuaciones diferenciales lineales homogéneas con coeficientes constantes reales.

\begin{multicols}{2}
\begin{enumerate}[itemsep=0.15cm]
    \item $y''' - 6y'' + 11y' - 6y = 0$
    \item $y''' - 3y'' + 3y' - y = 0$
    \item $y''' + y' = 0$
    \item $y^{(4)} - 16y = 0$
    \item $y^{(4)} + 8y'' + 16y = 0$
    \item $y''' - y'' - y' + y = 0$
    \item $y^{(4)} - 5y'' + 4y = 0$
    \item $y''' + 2y'' + 4y' + 8y = 0$
    \item $y^{(5)} - 2y^{(4)} + y''' = 0$
    \item $y^{(4)} + 2y'' + y = 0$
\end{enumerate}
\end{multicols}

\subsection*{B) Problemas de Valor Inicial (PVI)}
Resuelve los siguientes problemas de valor inicial para ecuaciones diferenciales lineales de orden superior.

\begin{enumerate}[itemsep=0.2cm]
    \item $y''' - 2y'' - y' + 2y = 0, \quad y(0) = 1, \quad y'(0) = 2, \quad y''(0) = 0$
    \item $y''' + 4y' = 0, \quad y(0) = 0, \quad y'(0) = 4, \quad y''(0) = 8$
    \item $y''' - 3y'' + 3y' - y = 0, \quad y(0) = 1, \quad y'(0) = 2, \quad y''(0) = 3$
    \item $y^{(4)} - y = 0, \quad y(0) = 1, \quad y'(0) = 0, \quad y''(0) = 1, \quad y'''(0) = 0$
    \item $y''' + y'' + y' + y = 0, \quad y(0) = 2, \quad y'(0) = -1, \quad y''(0) = 0$
\end{enumerate}

\subsection*{C) Reconstrucción de la Ecuación Diferencial}
Para cada uno de los siguientes conjuntos fundamentales de soluciones, reconstruye la ecuación diferencial lineal homogénea con coeficientes constantes reales de orden mínimo $P(D)y = 0$.

\begin{multicols}{2}
\begin{enumerate}[itemsep=0.15cm]
    \item $S = \{1, \, e^{2x}, \, e^{-2x}\}$
    \item $S = \{e^x, \, x e^x, \, x^2 e^x\}$
    \item $S = \{\cos(3x), \, \sin(3x), \, e^x\}$
    \item $S = \{1, \, x, \, e^{4x}\}$
    \item $S = \{e^{-x}, \, e^x \cos(2x), \, e^x \sin(2x)\}$
\end{enumerate}
\end{multicols}

\subsection*{D) Problemas Conceptuales de Deducción y Orden Mínimo}
Resuelve los siguientes problemas de deducción aplicando las propiedades de los polinomios característicos con coeficientes reales:

\begin{enumerate}[itemsep=0.25cm]
    \item \textbf{Solución compleja:} Una ecuación diferencial lineal homogénea con coeficientes constantes reales admite como solución particular $y_1(x) = e^{3x} \sin(2x)$.
    \begin{enumerate}[label=(\alph*), itemsep=0.08cm]
        \item ¿Cuál es el orden mínimo que puede tener dicha ecuación diferencial?
        \item Determina las raíces del polinomio característico y escribe la correspondiente ecuación diferencial.
        \item Si además se sabe que $y_2(x) = 5$ también es solución de la misma ecuación, ¿cuál es el nuevo orden mínimo y la ecuación diferencial resultante?
    \end{enumerate}

    \item \textbf{Raíz compleja repetida:} Una ecuación diferencial lineal homogénea con coeficientes reales constantes admite como solución $y_1(x) = x e^{-x} \cos(x)$.
    \begin{enumerate}[label=(\alph*), itemsep=0.08cm]
        \item Identifica todas las raíces características que esta solución implica, indicando sus multiplicidades.
        \item Determina el orden mínimo $n$ de la ecuación diferencial.
        \item Escribe la ecuación diferencial en notación de operadores $P(D)y = 0$.
    \end{enumerate}

    \item \textbf{Reconstrucción con condiciones iniciales:} Una ecuación diferencial lineal homogénea de tercer orden con coeficientes constantes reales tiene como soluciones $y_1(x) = \sin(x)$ y $y_2(x) = e^{2x}$.
    \begin{enumerate}[label=(\alph*), itemsep=0.08cm]
        \item Deduce el conjunto completo de las 3 raíces características y escribe la ecuación diferencial.
        \item Resuelve el problema de valor inicial sujeto a las condiciones $y(0) = 0$, $y'(0) = 3$ y $y''(0) = 5$.
    \end{enumerate}
\end{enumerate}

\end{document}
